\section{Introduction}
\label{sec:intro}

Six-player Canadian Fish is a team game of pure public information flow. Fifty-four cards are dealt
nine each into two alternating teams of three; a player asks a named opponent for a named card,
must already hold another card of that card's six-card half-suit, must not hold the card itself, and
keeps the turn on a hit and loses it on a miss; a team scores a half-suit by naming, correctly and
in full, which of its members holds each of the six cards. Every ask, every hit, every miss and
every declaration is seen by all six players. That single property makes the entire hidden state the
\emph{initial deal}, and it makes the posterior over that deal exactly computable --- which is why
this project has been able to treat inference as a solved problem and spend three studies on policy.

The three prior FishBots are, in order: v0.3, a fitted linear policy in a TypeScript simulator;
v0.4, a rules-exact C++ engine with an exact reference posterior, a fast deployed approximation, a
twenty-feature ask score and a value-based declaration rule; and v0.5, which diagnosed and removed a
failure mode in v0.4 --- a two-question deadlock in mirror play that ended in misdeclaration --- and
reported, in terms, that removing it was worth almost nothing in win rate.

That report is the starting point here. A project whose last two releases produced no measurable
strength gain is either at the ceiling of its policy class or is measuring badly, and the two
possibilities call for opposite responses. This study answers the question by building the
mechanisms the accumulated design register ranked highest, measuring each at a resolution the
project had not previously used, and then measuring the measurement.

\paragraph{Contributions.}
\begin{enumerate}
\item \textbf{Exchangeable ties are irreducible} (\S\ref{sec:diag-irreducible}). More than half of
  v0.5's ask decisions end in a bit-for-bit tie; the tied candidates are two cards of one half-suit
  at one target; the \emph{exact} posterior separates \vsixTieExactSepPctFive\% of them; and every
  tie-break rule realises the same hit rate. The largest item on the inherited register is a
  property of the information set, not a defect.
\item \textbf{Exactness is not the same as informativeness} (\S\ref{sec:belief-exactworse}). The
  exact posterior is a worse predictor of card location than the deployed approximation, because
  the approximation carries a soft policy prior and the exact object is exact under a uniform prior
  over deals. This closes an experiment three studies have listed as outstanding.
\item \textbf{The optimizer's curse in determinized search} (\S\ref{sec:search-curse}). A search
  that determinizes the deal but reconstructs all six continuation players at their own information
  sets --- not a double-dummy rollout --- is \vsixCurseGap{} points worse than its own blueprint
  when the action is chosen by an unguarded argmax, and a paired deviation rule recovers all of it.
  Guarded, it is the one mechanism in this study that beats a static rule: \vsixSearchLift\% against
  \vsix{} itself over \vsixSearchLiftN{} games with every cell above parity.
\item \textbf{A fitting harness that moves the policy} (\S\ref{sec:fitting}). Four measured defects
  in the previous optimiser, a repair, and a recovery test the previous one provably fails. The
  repaired optimiser finds a vector that trades nothing head to head for a large, replicated gain
  against the deception archetype.
\item \textbf{An evidence standard} (\S\ref{sec:protocol-paired}). Five mechanisms in this study
  cleared a $95\%$ paired interval at the per-cell budget this project has used since v0.3 and
  returned exactly zero at three times it, on two disjoint banks.
\item \textbf{Engine corrections} (\S\ref{sec:engine-repairs}, \S\ref{sec:corrections}), including
  an exact-inference object that returned another observer's tables whenever two observers held one,
  and a declaration pre-gate audit that had been dead code since v0.5.
\end{enumerate}

\paragraph{What this paper does not claim.} \vsix{} beats \vfive{} head to head by
\vsixBigFive\% $-$ 50, which is \emph{small}: about nine tenths of a point, resolved only at
\vsixBigFiveN{} games. It is not a large win-rate improvement, and at the per-cell budget this
project has always used it reads as a null. The substantive claims are robustness across playstyles,
lower measured exploitability, correctness, and a substantially narrowed design space.
